\documentclass[10pt,a4paper]{amsart}
\usepackage[margin=19mm]{geometry}
\usepackage{amsmath,amssymb,mathtools}
\usepackage[scr=rsfs,cal=euler]{mathalfa}
\usepackage{xfrac}
\usepackage[unicode,hidelinks,bookmarks=false]{hyperref}
\newcommand{\ie}{i.e.\ }
\newcommand{\cf}{cf.\ }
\newcommand{\resp}{respectively\ }
\newcommand{\Subsection}{\subsection}
\providecommand{\nobreakdash}{\nobreak\hbox{-}}
\setlength{\emergencystretch}{2em}
\multlinegap0pt
\usepackage{amssymb}
\newtheorem{theorem}{Theorem}
\title{\textbf{Fourier decay of equilibrium states and the Fibonacci Hamiltonian}}
\author{Ga\'etan Leclerc}
\date{Source: arXiv:2507.23731v2 (CC BY 4.0)}
\begin{document}
\maketitle

\noindent\emph{Source and licence.} \textbf{Fourier decay of equilibrium states and the Fibonacci Hamiltonian}. Version arXiv:2507.23731v2 (CC BY 4.0), distributed by arXiv under the Creative Commons Attribution 4.0 International licence, http://creativecommons.org/licenses/by/4.0/, as recorded in the arXiv licence metadata for this exact version and shown on the abstract page https://arxiv.org/abs/2507.23731v2. Full licence notice: \texttt{licenses/arxiv-2507.23731-CC-BY-4.0.txt}.\par\smallskip

\noindent\textit{Editorial scope.} Counted unit: the article's theorem th:IDS3 (source0.tex lines 263--266). The article's complete original proof of it is delivered with it (source0.tex lines 268--327, one \texttt{proof} environment of 60 lines). No mathematics is written, repaired or re-derived: the statement and the proof are the article's own lines, in the article's own notation, and no retained line is substituted or re-flowed.  Retained uncounted as the source-local context the proof needs: lines 105--108; lines 153--157; lines 208--212; lines 218--223.  Nothing was omitted from any retained span, so the package carries no omission record.  No retained line contains a citation, so the wrapper carries no bibliography and no bibliography entry.  No retained line contains a figure environment, an image-inclusion command or drawing code, so no image is delivered and the package declares no asset dependency.  Editorial formatting only, in addition: an AMS \texttt{amsart} preamble that restates the article's own declarations and macros used by the retained lines, namely the environments \texttt{theorem} on separate counters, together with the standard packages those lines need.  The retained environments print in this document as Theorem 1 in order of appearance, whereas the article prints the same items with the numbering of its own full document.  The complete native source of the article as distributed in the arXiv e-print for this exact version is bundled unchanged as \texttt{sources/182508/source0.tex}.
\begin{theorem}\label{th:IDS1}
There exists a closed and discrete subset $\mathcal{V}\subset [0,\infty)$ that doesn't contain $0$ such that the following holds. For any $V \in [0,\infty) \setminus \mathcal{V}$, there exists $\rho(V)>0$ and $C(V) \geq 1$ such that:
$$ \forall t \in \mathbb{R}, \ \Big| \int_\omega \langle \delta_0, e^{-it H_{V,\omega}} \delta_0 \rangle_{\ell^2(\mathbb{Z})} d\omega \Big| \leq \frac{C}{1+|t|^\rho}. $$
\end{theorem}

\begin{theorem}\label{th:IDS2}
Denote by $N_V \in \mathcal{P}(\Sigma_V)$ the density of states measure associated with the Fibonacci Hamiltonian. There exists a closed and discrete subset $\mathcal{V}\subset [0,\infty)$ that doesn't contain $0$ such that for any $V \in [0,\infty) \setminus \mathcal{V}$, and for any $\varepsilon>0$, we have:
$$ \underline{\text{dim}}_{F,C^{1+\varepsilon}}(N_V)>0 .$$
In particular, for any $C^{1+}$-diffeomorphism $\varphi:\mathbb{R} \rightarrow \mathbb{R}$, we have $\dim_F(\varphi(\Sigma_V)) > 0$.
\end{theorem}

\begin{theorem}\label{th:main1}
Let $f:M \rightarrow M$ be a $C^\infty$ Axiom A diffeomorphim on a smooth and complete riemannian surface $M$. Let $\Omega$ be a basic set for $f$. Suppose that $|\det df| = 1$ on $\Omega$, and suppose that $E^s \notin C^2(\Omega,TM)$ or $E^u \notin C^2(\Omega,TM)$. Then the measure of maximal entropy $\mu \in \mathcal{P}(\Omega)$ satisfies
$$ \underline{\text{dim}}_{F,C^{1+\alpha}}(\mu) > 0,$$
for any $\alpha>0$.
\end{theorem}

\begin{theorem}\label{th:main2}
Let $f:M \rightarrow M$ be a $C^\infty$ Axiom A diffeomorphim on a smooth and complete riemannian surface $M$. Let $\Omega$ be a basic set for $f$. Suppose that $|\det df| = 1$ on $\Omega$, and suppose that $E^s \notin C^2(\Omega,TM)$ or $E^u \notin C^2(\Omega,TM)$. Let $L \subset M$ be a line transverse to $W^s_{loc}$, and consider the associated stable holonomy $\pi_L : \Omega \rightarrow L$. Denote by $\mu \in \mathcal{P}(\Omega)$ the measure of maximal entropy, and $\nu := (\pi_L)_*(\mu) \in \mathcal{P}(L)$ the pushforward of $\mu$ by $\pi_L$. Then:
$$ \underline{\text{dim}}_{F,C^{1+\alpha}}( \nu ) > 0.$$
In particular, if we identify $\nu$ with a measure $\tilde{\nu}$ on the real line via some $C^{1+}$ parametrization of $L$, we find:
$$ \exists \rho >0, \exists C \geq 1, \forall \xi \in \mathbb{R}^*, \ \Big| \int_\mathbb{R} e^{i x \xi} d\tilde{\nu}(x) \Big| \leq C |\xi|^{-\rho}. $$
\end{theorem}

\begin{theorem}\label{th:IDS3}
Denote by $\mu_V \in \mathcal{P}(\Omega_V)$ the measure of maximal entropy for $(T_{|S_V},\Omega_V)$, and $N_V \in \mathcal{P}(\Sigma_V)$ the integrated density of state for the Fibonacci Hamiltonian. There exists a closed and discrete subset $\mathcal{V}\subset [0,\infty)$ that doesn't contain $0$ such that the following holds. For any $V \in [0,\infty) \setminus \mathcal{V}$, we have for all $\alpha \in (0,1)$:
$$   \underline{\text{dim}}_{F,C^{1+\alpha}}(\mu_V) > 0 \quad \text{and} \quad \underline{\text{dim}}_{F,C^{1+\alpha}}(N_V) > 0 . $$
\end{theorem}

\begin{proof}[Proof (of Theorem \ref{th:IDS1}, Theorem \ref{th:IDS2} and Theorem \ref{th:IDS3})] 
The goal of this proof is to compute an Anosov cocycle in some coordinates depending on $V$. We will show that this cocycle will be non-vanishing for small enough $V$. But then, by analyticity in $V$, it will be non-vanishing almost everywhere in $V$. Let us start computing.\\

Iterating once the Trace map, we find
$$ T^2(x,y,z) = ((4x^2-1)y-2xz,2xy-z,z). $$
The set of $2$-periodic points of $T$ is then given by $\text{Per}_2(T) = \{ (t,\frac{t}{2t-1},t), \ t \in \mathbb{R} \setminus \{ 1/2\} \}$. If $t = 1+\varepsilon > 1$ is close to one, we have $p_t := (t,\frac{t}{2t-1},t) \simeq (1,1,1) + \varepsilon(1,-1,1)$. In particular, the $y$-coordinate is less than one if $t>1$. \\

Denote by $t_V>1$ the only real number close to $1^+$ when $V$ is small such that $p_{t_V} \in S_V = \{ (x,y,z) \in \mathbb{R}^3, \  x^2+y^2+z^2 - 2xyz = 1 + V^2/4 \} $ for $V>0$. The relation 
$$ t_V^2 + \frac{t_V^2}{(2t_V-1)^2} +t_V^2 - 2 \frac{t_V^3}{2t_V-1} = 1+ \frac{V^2}{4} $$
can be rewritten as $$ \frac{(t_V-1)^2(4t_V^2+2t_V-1)}{(2t_V-1)^2} = \frac{V^2}{4}, $$
which reveals the useful estimate $$t_V = 1+\frac{V}{2\sqrt{5}} + O(V^2).$$
The point $p_V := p_{t_V} \in S_V$ is then a fixed point for the Axiom A, area-preserving, smooth map $T^2_{|S_V}$. Our goal now is to construct adapted coordinates centered at $p_V$, and to compute the Anosov cocycle in these coordinates.  \\

Let us begin by introducing a first parametrisation of the surface $S_V$, $V>0$. These coordinates are not going to be adapted yet but we will modify them later. Let $(x,y,z) \in \mathbb{R}^3$ with $y<1$, close to $(1,1,1)$. Then $(x,y,z) \in S_V$ iff $$x^2+y^2+z^2 - 2x y z = 1+\frac{V^2}{4}. $$
Solving in $y$ yields
$$ y = xz - \sqrt{(x^2-1)(z^2-1) + \frac{V^2}{4}} := y_V(x,z). $$
We can then parametrize our surface $S_V$ close to the fixed point $p_V$ by 
$(x,z) \mapsto (x,y_V(x,z),z)$.
Notice that in these coordinates, $p_V$ is represented by $(t_V,t_V)$. For future use, we will need to derive asymptotics in $V>0$ of derivatives of $y_V$ at $(t_V,t_V)$, using the relation $t_V = 1 + \frac{V}{2 \sqrt{5}} + O(V^2)$. \\

Notice that $y_V(x,z)=y_V(z,x)$, so that on the diagonal, $\partial_x y_V(t,t) = \partial_z y_V(t,t),$ $\partial_x^2 y_V(t,t) = \partial_z^2 y_V(t,t)$, etc. By construction, we know that $$y_V(t_V,t_V) = \frac{t_V}{2t_V-1} = 1-\frac{V}{2\sqrt{5}} + O(V^2).$$
Then, for the first-order derivative, we have
$$ \partial_x y_V(x,z) = z - \frac{x(z^2-1)}{\sqrt{(x^2-1)(z^2-1)+\frac{V^2}{4}}}, $$
which gives, recalling $t_V=1+\frac{V}{2\sqrt{5}} + O(V^2)$: $$ \partial_x y_V(t_V,t_V) = \partial_z y_V(t_V,t_V) = t_V - \frac{t_V(t_V^2-1)}{\sqrt{(t_V^2-1)^2+\frac{V^2}{4}}} = 1 - \frac{V/\sqrt{5}}{\sqrt{\frac{V^2}{5}+\frac{V^2}{4}}} + O(V) = \frac{1}{3} + O(V). $$
Similar direct computations reveals the behavior of $y_V$ up to the third order in $(x,z)$, that we summarize here:
\begin{align*}
    y_V(t_V,t_V) &= 1-\frac{V}{2\sqrt{5}} + O(V^2)  & \partial_x y_V(t_V,t_V) &= \frac{1}{3} + O(V) \\
    \partial_{xx}^2 y_V(t_V,t_V) &= \frac{8\sqrt{5}}{27} V^{-1} + O(1) &
    \partial_{xz}^2 y_V(t_V,t_V) &= -\frac{28\sqrt{5}}{27} V^{-1} + O(1) \\
    \partial_{xxz}^3 y_V(t_V,t_V) &= \frac{320}{81} V^{-2} + O(V^{-1}) &
    \partial_{xxx}^3 y_V(t_V,t_V) &= -\frac{160}{81} V^{-2} + O(V^{-1})
\end{align*}
We are now ready to express the dynamics in our coordinate system. Denoting $f_V(x,z)$ the map $T^2_{|S_V}$ in our coordinates near $p_V$, we find the formula
$$ f_V\begin{pmatrix} x \\ z \end{pmatrix} = \begin{pmatrix} (4x^2-1) \ y_V(x,z)-2xz \\ x \end{pmatrix} .$$
A direct computation allows us to find the derivatives of $f_V$ at $(t_V,t_V)$. We can summarize these informations as a Taylor expansion like so:
$$ f_V\Big( \begin{pmatrix} t_V \\ t_V \end{pmatrix} + \begin{pmatrix} x \\ z \end{pmatrix} \Big)  =  \begin{pmatrix} t_V \\ t_V \end{pmatrix}+\begin{pmatrix} 7 + O(V) & -1 + O(V) \\ 1 & 0 \end{pmatrix}  \begin{pmatrix} x \\ z \end{pmatrix}  $$
$$ + \begin{pmatrix} \frac{4\sqrt{5}}{9} x^2 -\frac{28\sqrt{5}}{9}  xz + \frac{4\sqrt{5}}{9} z^2 \\ 0 \end{pmatrix} (V^{-1} + O(1)) $$ $$ +\begin{pmatrix} -\frac{80}{81} x^3 + \frac{160}{27} x^2 z + \frac{160}{27} x z^2 - \frac{80}{81} z^3 \\ 0 \end{pmatrix} (V^{-2} + O(V^{-1}))  + O_V(x^4+z^4). $$
For example, one can compute the Jacobian matrix at $(t_V,t_V)$ as follow:
$$ J_{f_V}(x,z) = \begin{pmatrix} 8x \cdot y_V(x,z) + (4x^2-1) \partial_xy_V(x,z) -2z & (4x^2-1) \partial_z y_V(x,z) - 2x \\ 1 & 0 \end{pmatrix} $$ 
and so
$$ J_{f_V}(t_V,t_V) = \begin{pmatrix} 8t_V \cdot y_V(t_V,t_V) + (4t_V^2-1)\partial_x y_V(t_V,t_V) - 2 t_V & (4t_V^2-1) \partial_z y_V(t_V,t_V) - 2t_V \\ 1 & 0 \end{pmatrix} $$ $$= \begin{pmatrix} 7  + O(V) & -1  + O(V) \\ 1 & 0 \end{pmatrix} . $$
Higher derivatives are computed similarly. Notice that $J_V := J_{f_V}(t_V,t_V)$ is analytic in $V$ and can be extended on a neighborhood of $V=0$. Notice also that we have
$$ P_0^{-1} J_0 P_0 = D_0, $$
where $$ D_0 := \begin{pmatrix} \frac{7+3\sqrt{5}}{2} & 0 \\ 0 & \frac{7-3\sqrt{5}}{2}  \end{pmatrix}  \quad , \quad P_0 := \frac{1}{2}\begin{pmatrix} 7+3\sqrt{5} & 7-3\sqrt{5} \\ 2 & 2 \end{pmatrix} \quad , \quad P_0^{-1} = \frac{1}{30}\begin{pmatrix} 2\sqrt{5} & 15-7 \sqrt{5} \\ -2\sqrt{5} & 15+7\sqrt{5} \end{pmatrix}.$$
Notice how $J_0$ has distinct real eigenvalues: by analyticity of $V \mapsto J_V$, the spectrum of $J_V$ is still real and distinct for small $V>0$. This in fact holds true for any $V>0$ since we know that $T_{|\Omega_V}^2$ is always hyperbolic. Moreover, since $T^2_{|S_V}$ is area preserving, we know that $\det J_V=1$. We can then find an analytic family of matrices $P_V,D_V$ such that, for all $V>0$,
$$ P_V^{-1} J_V P_V = D_V $$
where $D_V$ is diagonal (with coefficients $\lambda_V,\mu_V$). By analyticity, we then have $D_V = D_0 + O(V)$, $P_V=P_0 + O(V)$, and $P_V^{-1} = P_0^{-1} + O(V)$.
Now, we are ready to introduce our adapted coordinates: we let
$$ \tilde{f}_V(X,Y) := P_V^{-1} f_V\Big( \begin{pmatrix} t_V \\ t_V \end{pmatrix} + P_V \begin{pmatrix} X \\ Y \end{pmatrix} \Big). $$
These coordinates are adapted. To estimate the derivatives $\partial^\alpha \tilde{f}_V(0,0)$, one can proceed as follow. Let $(h_i)$ be some vector of the form $(1,0)$ or $(0,1)$. The $k$-th differential can be written as
$$ (d^k \tilde{f}_V)_{(0,0)}(h_1,\dots,h_k) = P_V^{-1} (d^{k}f_V)_{(t_V,t_V)}\Big( P_V h_1, \dots , P_V h_k \Big) $$
$$ = P_0^{-1} (d^{k}f_V)_{(t_V,t_V)}\Big( P_0 h_1, \dots , P_0 h_k \Big) + O(V \cdot \|(d^k f_V)_{(t_V,t_V)}\| ). $$
To find the derivatives of $\tilde{f}_V$, we then compute the Taylor coefficients of \\ $P_0^{-1} f_V\Big( \begin{pmatrix} t_V \\ t_V \end{pmatrix} + P_0 \begin{pmatrix} X \\ Y \end{pmatrix} \Big)$, by using the known Taylor estimate on $f_V$ at $(t_V,t_V)$, replacing the variables $(x,z)$ with $(X,Y)$. A direct computation yields:
$$ \tilde{f}_V(X,Y) = \begin{pmatrix} \frac{7+3\sqrt{5}}{2} + O(V) & 0 \\ 0 &  \frac{7-3\sqrt{5}}{2} + O(V) \end{pmatrix} \begin{pmatrix} X \\ Y \end{pmatrix} + \frac{20}{3} XY \begin{pmatrix} -1 \\ 1 \end{pmatrix} (V^{-1} + O(1)) + (X^2+Y^2) O(1) $$ $$+ \frac{40}{9} X Y \big( (5+3 \sqrt{5})X - (5-3 \sqrt{5})Y \big) \begin{pmatrix} 1 \\ -1 \end{pmatrix} (V^{-2} + O(V^{-1})) + (X^3+Y^3)O(V^{-1}) + O_V(X^4+Y^4) .$$
We can now estimate the Anosov cocycle in these coordinates. Denoting $\tilde{f}_V(X,Y) =  \begin{pmatrix} F_V(X,Y) \\ G_V(X,Y) \end{pmatrix}$, we have: $\mathcal{A}_V :=$ 
$$ \frac{\partial^3_{XYY} G_V(0,0)}{\mu_V} - \frac{\partial_{XY}^2 F_V(0,0) \partial_{XY}^2 G_V(0,0)}{\lambda_V-1} - \frac{\partial_{XX}^2 G_V(0,0) \partial_{YY}^2 F_V(0,0)}{1-\mu_V^3} - \frac{\partial_{XY}^2 G_V(0,0) \partial_{YY}^2 G_V(0,0)}{\mu_V^2}  $$
$$ = \frac{\partial^3_{XYY} G_V(0,0)}{\mu_0} - \frac{\partial_{XY}^2 F_V(0,0) \partial_{XY}^2 G_V(0,0)}{\lambda_0-1} + O(V^{-1}) $$
$$ =  \frac{160}{9}\frac{5-3\sqrt{5}}{7-3\sqrt{5}} V^{-2} - \frac{40}{9} \frac{2}{5+3\sqrt{5}} V^{-2} + O(V^{-1}) $$
$$ = - \frac{140+76\sqrt{5}}{3} \ V^{-2} + O(V^{-1}), $$
which is nonvanishing for $V>0$ small enough. Notice finally that the Anosov cocycle $\mathcal{A}_V$ is analytic in $V$ on the domain $V>0$, and doesn't vanish on a neighborhood of $V=0$. It follows that it only vanishes at most on a discrete subset $\mathcal{V} \subset [0,\infty)$ that doesn't contains $0$. It follows that $E^s \text{ or } E^u \notin C^2$ on $\Omega_V$ for $V \notin \mathcal{V}$, and so Theorem \ref{th:main1} and Theorem \ref{th:main2} apply. \end{proof}

\end{document}
